\documentclass[10pt,a4paper]{amsart}
\usepackage[margin=19mm]{geometry}
\usepackage{amsmath,amssymb,mathtools}
\usepackage[scr=rsfs,cal=euler]{mathalfa}
\usepackage{xfrac}
\usepackage[unicode,hidelinks,bookmarks=false]{hyperref}
\newcommand{\ie}{i.e.\ }
\newcommand{\cf}{cf.\ }
\newcommand{\resp}{respectively\ }
\newcommand{\Subsection}{\subsection}
\providecommand{\nobreakdash}{\nobreak\hbox{-}}
\setlength{\emergencystretch}{2em}
\multlinegap0pt
\usepackage{amssymb}
\usepackage{float}
\usepackage{xcolor}
\usepackage{algorithm}\usepackage{algpseudocode}
\newtheorem{proposition}{Proposition}
\title{Scalable $s$--step Preconditioned Conjugate Gradient with Chebyshev Basis and Gauss--Seidel Gram Solve}
\author{Pasqua D'Ambra, Massimo Bernaschi, Mauro G. Carrozzo, Stephen Thomas}
\date{Source: arXiv:2603.09790v3 (CC BY 4.0)}
\begin{document}
\maketitle

\noindent\emph{Source and licence.} Scalable $s$--step Preconditioned Conjugate Gradient with Chebyshev Basis and Gauss--Seidel Gram Solve. Version arXiv:2603.09790v3 (CC BY 4.0), distributed by arXiv under the Creative Commons Attribution 4.0 International licence, http://creativecommons.org/licenses/by/4.0/, as recorded in the arXiv licence metadata for this exact version and shown on the abstract page https://arxiv.org/abs/2603.09790v3. Full licence notice: \texttt{licenses/arxiv-2603.09790-CC-BY-4.0.txt}.\par\smallskip

\noindent\textit{Editorial scope.} Counted unit: the article's proposition prop:fgs-residual (source0.tex lines 622--664). The article's complete original proof of it is delivered with it (source0.tex lines 665--762, one \texttt{proof} environment of 98 lines). No mathematics is written, repaired or re-derived: the statement and the proof are the article's own lines, in the article's own notation, and no retained line is substituted or re-flowed.  Retained uncounted as the source-local context the proof needs: lines 597--602.  Nothing was omitted from any retained span, so the package carries no omission record.  No retained line contains a citation, so the wrapper carries no bibliography and no bibliography entry.  No retained line contains a figure environment, an image-inclusion command or drawing code, so no image is delivered and the package declares no asset dependency.  Editorial formatting only, in addition: an AMS \texttt{amsart} preamble that restates the article's own declarations and macros used by the retained lines, namely the environments \texttt{proposition} on separate counters, together with the standard packages those lines need.  The retained environments print in this document as Proposition 1 in order of appearance, whereas the article prints the same items with the numbering of its own full document.  The complete native source of the article as distributed in the arXiv e-print for this exact version is bundled unchanged as \texttt{sources/196088/source0.tex}.
\begin{equation}
\label{eq:fgs-iteration}
(I+L)\widetilde{\alpha}^{(\nu+1)}
=
\widetilde{m}-L^\top\widetilde{\alpha}^{(\nu)}.
\end{equation}

\begin{proposition}[Residual after $\nu$ FGS sweeps]
\label{prop:fgs-residual}
Let
\[
W_c=I+L+L^\top,
\]
where $L$ is strictly lower triangular. Let
$\widetilde{\alpha}^{(0)}=0$, and let
$\widetilde{\alpha}^{(\nu)}$ be obtained after $\nu$ forward
Gauss--Seidel sweeps applied to
\[
W_c\widetilde{\alpha}=\widetilde{m}.
\]
Then the residual of the normalized system satisfies
\begin{equation}
\label{eq:fgs-residual-formula}
\widetilde{\rho}^{(\nu)}
:=
\widetilde{m}
-
W_c\widetilde{\alpha}^{(\nu)}
=
(-1)^\nu
\left(
L^\top(I+L)^{-1}
\right)^\nu
\widetilde{m}.
\end{equation}
Consequently,
\begin{equation}
\label{eq:fgs-residual-bound}
\left\|
\widetilde{\rho}^{(\nu)}
\right\|
\leq
\left\|
L^\top(I+L)^{-1}
\right\|^\nu
\left\|
\widetilde{m}
\right\|.
\end{equation}
\end{proposition}

\begin{proof}
The proof proceeds by induction. For $\nu=1$, the FGS iteration with
$\widetilde{\alpha}^{(0)}=0$ gives
\[
(I+L)\widetilde{\alpha}^{(1)}=\widetilde{m},
\]
and hence
\[
\widetilde{\alpha}^{(1)}
=
(I+L)^{-1}\widetilde{m}.
\]
Therefore,
\begin{align*}
\widetilde{\rho}^{(1)}
&=
\widetilde{m}
-
W_c\widetilde{\alpha}^{(1)}
\\
&=
\widetilde{m}
-
(I+L+L^\top)(I+L)^{-1}\widetilde{m}
\\
&=
-L^\top(I+L)^{-1}\widetilde{m},
\end{align*}
which proves~\eqref{eq:fgs-residual-formula} for $\nu=1$.
Assume that~\eqref{eq:fgs-residual-formula} holds for some $\nu\geq1$.
From the FGS update~\eqref{eq:fgs-iteration},
\[
(I+L)\widetilde{\alpha}^{(\nu+1)}
=
\widetilde{m}
-
L^\top\widetilde{\alpha}^{(\nu)}.
\]
Using
\[
(I+L+L^\top)(I+L)^{-1}
=
I+L^\top(I+L)^{-1},
\]
we obtain
\begin{align*}
\widetilde{\rho}^{(\nu+1)}
&=
\widetilde{m}
-
W_c\widetilde{\alpha}^{(\nu+1)}
\\
&=
\widetilde{m}
-
(I+L+L^\top)(I+L)^{-1}
\left(
\widetilde{m}
-
L^\top\widetilde{\alpha}^{(\nu)}
\right)
\\
&=
-L^\top(I+L)^{-1}
\left(
\widetilde{m}
-
(I+L+L^\top)\widetilde{\alpha}^{(\nu)}
\right)
\\
&=
-L^\top(I+L)^{-1}
\widetilde{\rho}^{(\nu)}.
\end{align*}
By the induction hypothesis,
\[
\widetilde{\rho}^{(\nu)}
=
(-1)^\nu
\left(
L^\top(I+L)^{-1}
\right)^\nu
\widetilde{m},
\]
and therefore
\[
\widetilde{\rho}^{(\nu+1)}
=
(-1)^{\nu+1}
\left(
L^\top(I+L)^{-1}
\right)^{\nu+1}
\widetilde{m}.
\]
This completes the induction. The bound
in~\eqref{eq:fgs-residual-bound} follows from the sub-multiplicativity
of the norm.
\end{proof}

\end{document}
